\chapter{Congruenze e anelli $\mbZ_m$}\label{cap:congruenze}

Nel Capitolo~\ref{cap:equivalenze} abbiamo visto la congruenza modulo $m$ come relazione di equivalenza. Qui ne studiamo le classi,
le \textbf{classi di resto}, e scopriamo che con esse si possono fare somme e prodotti: si ottengono gli anelli $\mbZ_m$, l'aritmetica
«dell'orologio». Per farlo introduciamo un'idea generale: le \textbf{congruenze} di una struttura algebrica, cioè le relazioni di
equivalenza che rispettano le operazioni, e le \textbf{strutture quoziente}.

% ──────────────────────────────────────────────────────────────
\section{Classi di resto}\label{sec:classi-resto}

Ricordiamo (Definizione~\ref{def:congruenza-modulo}) che $a \equiv_m b$ se $b - a$ è un multiplo di $m$, cioè se $m \mid b - a$.

\dfn[zm]{Classi di resto, l'insieme $\mbZ_m$}{
    Sia $m \in \mbZ$. Poniamo $\mbZ_m := \mbZ/{\equiv_m}$ e, per $a \in \mbZ$, indichiamo la classe di $a$ con
    \[ [a]_m := [a]_{\equiv_m} = \{a + mk \mid k \in \mbZ\} = a + m\mbZ. \]
    Le classi $[a]_m$ si chiamano \textbf{classi di resto} modulo $m$.
}

\dfn[mod]{L'operazione $\bmod$}{
    Per $a \in \mbZ$ e $m \in \mbZ \setminus \{0\}$ poniamo $a \bmod m := \min([a]_m \cap \mbN)$ (si scrive anche $a \,\%\, m$): il più piccolo numero
    naturale congruo ad $a$ modulo $m$.
}

\begin{concetto}
\begin{Proposizione}{$a \bmod m$ è il resto della divisione}{mod-resto}
    \begin{prerequisiti}
        \dip{Definizione}{def:mod}
        \dip{Teorema}{teo:divisione-euclidea}
    \end{prerequisiti}
    Siano $a \in \mbZ$, $m \neq 0$ e $a = mq + r$ con $0 \leq r < |m|$. Allora $a \bmod m = r$. In particolare $0 \leq a \bmod m < |m|$.
\end{Proposizione}
\begin{dimostrazione}
    $r = a - mq \in [a]_m \cap \mbN$. Se fosse $x \in [a]_m \cap \mbN$ con $x < r$, allora $x \equiv_m r$ e $0 < r - x < |m|$: $r - x$ sarebbe un multiplo di $m$
    non nullo di valore assoluto minore di $|m|$, impossibile. Quindi $r$ è il minimo.
\end{dimostrazione}
\end{concetto}

\begin{spiegazione}
    Per trovare il rappresentante «canonico» di una classe basta dividere: $101 \bmod 10 = 1$, \ $30093 \bmod 3 = 0$,
    \ $-7 \bmod 3 = 2$ (perché $-7 = 3 \cdot (-3) + 2$). Su una calcolatrice o in un linguaggio di programmazione l'operatore $\%$ può
    dare risultati negativi per $a < 0$: con la nostra definizione il risultato è sempre tra $0$ e $|m| - 1$.
\end{spiegazione}

\begin{concetto}
\begin{Teorema}{Descrizione esplicita di $\mbZ_m$}{zm-esplicito}
    \begin{prerequisiti}
        \dip{Definizione}{def:zm}
        \dip{Teorema}{teo:divisione-euclidea}
        \dip{Teorema}{teo:proprieta-classi}
    \end{prerequisiti}
    Sia $m \in \mbZ \setminus \{0\}$. Allora $\mbZ_m = \{[0]_m, [1]_m, \dots, [|m| - 1]_m\}$ e queste $|m|$ classi sono a due a due distinte; in particolare
    $|\mbZ_m| = |m|$.
\end{Teorema}
\begin{dimostrazione}
    Poiché $\equiv_m$ e $\equiv_{-m}$ coincidono, possiamo supporre $m > 0$.
    \begin{caso}{Ogni classe è una delle $[r]_m$}
        Sia $a \in \mbZ$ e sia $a = mq + r$ con $0 \leq r < m$. Allora $a - r = mq$, cioè $a \equiv_m r$, e quindi $[a]_m = [r]_m$.
    \end{caso}
    \begin{caso}{Le $[r]_m$ sono distinte}
        Siano $0 \leq i \leq j \leq m - 1$ con $[i]_m = [j]_m$. Allora $m \mid j - i$ con $0 \leq j - i < m$, quindi $j - i = 0$: $i = j$.
    \end{caso}
\end{dimostrazione}
\end{concetto}

\begin{spiegazione}
    Gli interi si «avvolgono» attorno a un orologio con $m$ posizioni: ogni posizione è una classe di resto, e i numeri che ci
    finiscono sopra differiscono per giri completi, cioè per multipli di $m$. Per $m = 5$:
    \begin{center}
    \begin{tikzpicture}
        \draw[black!40, line width=0.8pt] (0,0) circle (1.3);
        \foreach \r/\ang/\col in {0/90/classeA, 1/18/classeB, 2/-54/classeC, 3/-126/classeD, 4/162/classeE} {
            \node[circle, fill=\col, draw=black!40, minimum size=8mm, inner sep=0pt, font=\small] at (\ang:1.3) {$[\r]_5$};
        }
        \node[font=\scriptsize, anchor=south] at (90:1.75) {$\dots, -5, 0, 5, 10, \dots$};
        \node[font=\scriptsize, anchor=west] at (18:1.75) {$\dots, -4, 1, 6, 11, \dots$};
        \node[font=\scriptsize, anchor=west] at (-54:1.75) {$\dots, -3, 2, 7, 12, \dots$};
        \node[font=\scriptsize, anchor=east] at (-126:1.75) {$\dots, -2, 3, 8, 13, \dots$};
        \node[font=\scriptsize, anchor=east] at (162:1.75) {$\dots, -1, 4, 9, 14, \dots$};
    \end{tikzpicture}
    \end{center}
    La dimostrazione dice due cose: ogni intero cade in una delle $5$ posizioni (quella del suo resto), e due posizioni diverse
    non coincidono (la distanza tra $0, 1, \dots, 4$ è sempre minore di un giro intero).
\end{spiegazione}

% ──────────────────────────────────────────────────────────────
\section{Congruenze in una struttura algebrica}\label{sec:congruenze-strutture}

Vorremmo sommare e moltiplicare le classi di resto ponendo $[a]_m + [b]_m := [a + b]_m$. Ma una classe ha tanti rappresentanti: la
definizione ha senso solo se il risultato \textbf{non dipende da quali rappresentanti scegliamo}. Questo succede quando la
relazione di equivalenza «va d'accordo» con l'operazione.

\dfn[congruenza]{Compatibilità, congruenza}{
    Sia $s \neq \varnothing$ e sia $*$ un'operazione binaria interna su $s$. Una relazione di equivalenza $\sim$ su $s$ si dice
    \begin{itemize}
        \item \textbf{compatibile a sinistra} con $*$ se $\forall a, b, c \in s\,(a \sim b \rightarrow c * a \sim c * b)$;
        \item \textbf{compatibile a destra} con $*$ se $\forall a, b, c \in s\,(a \sim b \rightarrow a * c \sim b * c)$.
    \end{itemize}
    Se $*_1, \dots, *_n$ sono operazioni binarie interne su $s$, $\sim$ si dice una \textbf{congruenza} in $(s, *_1, \dots, *_n)$ se, per ogni $i$,
    \[ \forall a, b, c, d \in s\,\big((a \sim b \wedge c \sim d) \rightarrow a *_i c \sim b *_i d\big). \]
}

\begin{concetto}
\begin{Teorema}{Congruenze e compatibilità}{congruenza-compatibile}
    \begin{prerequisiti}
        \dip{Definizione}{def:congruenza}
        \dip{Definizione}{def:tipi-relazioni}
    \end{prerequisiti}
    Una relazione di equivalenza $\sim$ su $s$ è una congruenza in $(s, *_1, \dots, *_n)$ se e solo se è compatibile a destra e a sinistra
    con ogni operazione $*_i$.
\end{Teorema}
\begin{dimostrazione}
    Basta considerare una sola operazione $*$.
    \begin{caso}{$(\rightarrow)$}
        Sia $a \sim b$. Applicando la definizione di congruenza con $c \sim c$ (riflessività) si ottiene $a * c \sim b * c$: compatibilità a destra.
        Applicandola alle coppie $c \sim c$ e $a \sim b$ si ottiene $c * a \sim c * b$: compatibilità a sinistra.
    \end{caso}
    \begin{caso}{$(\leftarrow)$}
        Siano $a \sim b$ e $c \sim d$. Per la compatibilità a destra $a * c \sim b * c$; per quella a sinistra $b * c \sim b * d$. Per la
        transitività $a * c \sim b * d$.
    \end{caso}
\end{dimostrazione}
\end{concetto}

\begin{spiegazione}
    Per sostituire \textbf{entrambi} i fattori con elementi equivalenti, li si sostituisce \textbf{uno alla volta}: prima quello a sinistra
    (compatibilità a destra), poi quello a destra (compatibilità a sinistra), e si collegano i due passi con la transitività:
    \[ a * c \ \overset{\text{cambio } a \to b}{\sim}\ b * c \ \overset{\text{cambio } c \to d}{\sim}\ b * d. \]
\end{spiegazione}

\begin{concetto}
\begin{Proposizione}{Le operazioni quoziente}{operazioni-quoziente}
    \begin{prerequisiti}
        \dip{Definizione}{def:congruenza}
        \dip{Teorema}{teo:proprieta-classi}
        \dip{Definizione}{def:funzione}
    \end{prerequisiti}
    Sia $\sim$ una congruenza in $(s, *_1, \dots, *_n)$. Allora per ogni $i$ è ben definita l'operazione binaria interna su $s/{\sim}$
    \[ (*_i)_\sim \colon ([x]_\sim, [y]_\sim) \in s/{\sim} \times s/{\sim} \mapsto [x *_i y]_\sim \in s/{\sim}, \]
    detta \textbf{operazione quoziente}. La struttura $(s/{\sim}, (*_1)_\sim, \dots, (*_n)_\sim)$ si chiama \textbf{struttura quoziente}.
\end{Proposizione}
\begin{dimostrazione}
    Bisogna verificare che il risultato non dipenda dai rappresentanti. Se $[x] = [x']$ e $[y] = [y']$, allora $x \sim x'$ e $y \sim y'$
    (Teorema~\ref{teo:proprieta-classi}); essendo $\sim$ una congruenza, $x *_i y \sim x' *_i y'$, cioè $[x *_i y] = [x' *_i y']$.
\end{dimostrazione}
\end{concetto}

\begin{spiegazione}
    Si opera sulle classi \textbf{scegliendo un rappresentante qualunque per ciascuna}, facendo il conto con i rappresentanti e
    prendendo la classe del risultato. Ad esempio in $\mbZ_5$: $[3]_5 + [4]_5 = [7]_5 = [2]_5$; scegliendo altri rappresentanti,
    $[8]_5 + [-1]_5 = [7]_5$ ancora, e $[13]_5 + [9]_5 = [22]_5 = [2]_5$. La proposizione garantisce che vada sempre così.
\end{spiegazione}

\begin{esempio}{Una relazione di equivalenza che non è una congruenza}
    Su $\mbZ$ consideriamo l'operazione $a * b = a + b$ se $b$ è dispari e $a * b = a + \frac{b}{2}$ se $b$ è pari. La congruenza $\equiv_2$
    \textbf{non} è compatibile con $*$: $0 \equiv_2 2$, ma $1 * 0 = 1$ e $1 * 2 = 2$, che non sono congrui modulo $2$. Quindi «$[1]_2 * [0]_2$» non si può
    definire: dipenderebbe dal rappresentante scelto per $[0]_2$. Invece con l'operazione $a * b = 2ab$ la relazione $\equiv_2$ è una congruenza:
    il risultato è sempre pari, quindi tutti i risultati sono congrui tra loro.
\end{esempio}

\begin{concetto}
\begin{Teorema}{L'epimorfismo canonico}{epimorfismo-canonico}
    \begin{prerequisiti}
        \dip{Proposizione}{prop:operazioni-quoziente}
        \dip{Definizione}{def:omomorfismo}
        \dip{Teorema}{teo:epimorfismi}
        \dip{Proposizione}{prop:proiezione}
    \end{prerequisiti}
    Sia $\sim$ una congruenza in $(s, *_1, \dots, *_n)$. La proiezione canonica $\pi \colon x \in s \mapsto [x]_\sim \in s/{\sim}$ è un epimorfismo (per
    ogni operazione), detto \textbf{epimorfismo canonico}. Di conseguenza la struttura quoziente eredita associatività, commutatività,
    elementi neutri e simmetrici, e la distributività di un'operazione rispetto a un'altra. In particolare i quozienti di semigruppi,
    monoidi, gruppi (abeliani) e anelli (commutativi, unitari) sono strutture dello stesso tipo.
\end{Teorema}
\begin{dimostrazione}
    $\pi$ è suriettiva (Proposizione~\ref{prop:proiezione}) e, per definizione di operazione quoziente, $\pi(x *_i y) = [x *_i y] = [x] (*_i)_\sim [y]
    = \pi(x) (*_i)_\sim \pi(y)$: è un epimorfismo. Il Teorema~\ref{teo:epimorfismi} dà allora associatività, commutatività, neutri e simmetrici.
    Per la distributività: $[x]([y] + [z]) = [x(y + z)] = [xy + xz] = [x][y] + [x][z]$ (scrivendo $+$ e $\cdot$ per due delle operazioni).
\end{dimostrazione}
\end{concetto}

\begin{spiegazione}
    Ogni proprietà espressa da un'uguaglianza tra conti «passa al quoziente»: per verificarla sulle classi la si verifica sui
    rappresentanti e poi si prende la classe. Ad esempio la commutatività in $\mbZ_m$: $[a][b] = [ab] = [ba] = [b][a]$. Il neutro di $\mbZ_m$
    per la somma è $[0]$, l'opposto di $[a]$ è $[-a]$, il neutro per il prodotto è $[1]$.
\end{spiegazione}

% ──────────────────────────────────────────────────────────────
\section{Gli anelli $\mbZ_m$}\label{sec:anelli-zm}

\begin{concetto}
\begin{Proposizione}{La congruenza modulo $m$ è una congruenza di $(\mbZ, +, \cdot)$}{congruenza-z}
    \begin{prerequisiti}
        \dip{Teorema}{teo:congruenza-compatibile}
        \dip{Teorema}{teo:epimorfismo-canonico}
        \dip{Proposizione}{prop:prop-divisibilita}
    \end{prerequisiti}
    Per ogni $m \in \mbZ$, $\equiv_m$ è una congruenza in $(\mbZ, +, \cdot)$. Quindi $(\mbZ_m, +, \cdot)$, con
    \[ [a]_m + [b]_m := [a + b]_m, \qquad [a]_m \cdot [b]_m := [ab]_m, \]
    è un anello commutativo unitario, con zero $[0]_m$ e unità $[1]_m$.
\end{Proposizione}
\begin{dimostrazione}
    Poiché somma e prodotto in $\mbZ$ sono commutativi, basta la compatibilità a destra (Teorema~\ref{teo:congruenza-compatibile}). Sia
    $a \equiv_m b$, cioè $m \mid b - a$. Allora $(b + c) - (a + c) = b - a$ e $bc - ac = (b - a)c$ sono multipli di $m$: $a + c \equiv_m b + c$ e
    $ac \equiv_m bc$. La seconda parte segue dal Teorema~\ref{teo:epimorfismo-canonico}, perché $(\mbZ, +, \cdot)$ è un anello commutativo unitario.
\end{dimostrazione}
\end{concetto}

\begin{spiegazione}
    È l'aritmetica dell'orologio: in $\mbZ_{12}$, $[5] + [9] = [14] = [2]$ dice che se sono le $5$, tra $9$ ore saranno le $2$. Ecco le tavole
    di Cayley di $\mbZ_4$ (rappresentanti $0, 1, 2, 3$; scriviamo $r$ per $[r]_4$):
    \begin{center}
    \renewcommand{\arraystretch}{1.2}
    \begin{tabular}{c|cccc}
        $+$ & $0$ & $1$ & $2$ & $3$ \\ \hline
        $0$ & $0$ & $1$ & $2$ & $3$ \\
        $1$ & $1$ & $2$ & $3$ & $0$ \\
        $2$ & $2$ & $3$ & $0$ & $1$ \\
        $3$ & $3$ & $0$ & $1$ & $2$ \\
    \end{tabular}
    \qquad\qquad
    \begin{tabular}{c|cccc}
        $\cdot$ & $0$ & $1$ & $2$ & $3$ \\ \hline
        $0$ & $0$ & $0$ & $0$ & $0$ \\
        $1$ & $0$ & $1$ & $2$ & $3$ \\
        $2$ & $0$ & $2$ & \cellcolor{cherryred!20}$0$ & $2$ \\
        $3$ & $0$ & $3$ & $2$ & \cellcolor{seagreen!25}$1$ \\
    \end{tabular}
    \end{center}
    Si legge che $[2]_4 \cdot [2]_4 = [0]_4$: in $\mbZ_4$ c'è un \textbf{divisore dello zero}, quindi $\mbZ_4$ non è un dominio di integrità; invece
    $[3]_4 \cdot [3]_4 = [1]_4$: $[3]_4$ è invertibile.
\end{spiegazione}

\begin{concetto}
\begin{Teorema}{Quando $\mbZ_m$ è un campo}{zm-campo}
    \begin{prerequisiti}
        \dip{Proposizione}{prop:congruenza-z}
        \dip{Proposizione}{prop:campo-dominio}
        \dip{Definizione}{def:primo}
        \dip{Teorema}{teo:primi-irriducibili}
        \dip{Teorema}{teo:bezout}
        \dip{Teorema}{teo:zm-esplicito}
    \end{prerequisiti}
    Sia $m \in \mbZ \setminus \{0\}$. Sono equivalenti:
    \begin{enumerate}
        \item $(\mbZ_m, +, \cdot)$ è un campo;
        \item $(\mbZ_m, +, \cdot)$ è un dominio di integrità;
        \item $m$ è primo.
    \end{enumerate}
\end{Teorema}
\begin{dimostrazione}
    \begin{caso}{(1) $\rightarrow$ (2)}
        Ogni campo è un dominio di integrità (Proposizione~\ref{prop:campo-dominio}).
    \end{caso}
    \begin{caso}{(2) $\rightarrow$ (3)}
        In un dominio di integrità $[0] \neq [1]$, quindi $|m| = |\mbZ_m| \geq 2$, cioè $m \notin \{-1, 0, 1\}$. Sia $m \mid ab$: allora
        $[a][b] = [ab] = [0]$ e, per la legge di annullamento del prodotto, $[a] = [0]$ oppure $[b] = [0]$, cioè $m \mid a$ oppure $m \mid b$. Quindi $m$ è primo.
    \end{caso}
    \begin{caso}{(3) $\rightarrow$ (1)}
        $\mbZ_m$ è commutativo e unitario, con $[0] \neq [1]$ perché $|m| \geq 2$. Sia $[n] \neq [0]$, con $1 \leq n < |m|$ (Teorema~\ref{teo:zm-esplicito}). Poiché $m$
        è primo, è irriducibile: i divisori comuni di $m$ e $n$ sono tra $\pm 1, \pm m$, e $\pm m$ non dividono $n$. Quindi $m$ e $n$ sono coprimi e per
        Bézout $1 = nu + mv$ per certi $u, v$. Allora $[n][u] = [nu] = [1 - mv] = [1]$: ogni classe non nulla è invertibile.
    \end{caso}
\end{dimostrazione}
\end{concetto}

\begin{spiegazione}
    Confrontiamo le righe del prodotto (senza lo zero) in $\mbZ_6$ e in $\mbZ_7$:
    \begin{center}
    \renewcommand{\arraystretch}{1.15}
    \begin{tabular}{c|ccccc}
        $\mbZ_6$ & $1$ & $2$ & $3$ & $4$ & $5$ \\ \hline
        $1$ & $1$ & $2$ & $3$ & $4$ & $5$ \\
        $2$ & $2$ & $4$ & \cellcolor{cherryred!20}$0$ & $2$ & $4$ \\
        $3$ & $3$ & \cellcolor{cherryred!20}$0$ & $3$ & \cellcolor{cherryred!20}$0$ & $3$ \\
        $4$ & $4$ & $2$ & \cellcolor{cherryred!20}$0$ & $4$ & $2$ \\
        $5$ & $5$ & $4$ & $3$ & $2$ & \cellcolor{seagreen!25}$1$ \\
    \end{tabular}
    \qquad\qquad
    \begin{tabular}{c|cccccc}
        $\mbZ_7$ & $1$ & $2$ & $3$ & $4$ & $5$ & $6$ \\ \hline
        $1$ & \cellcolor{seagreen!25}$1$ & $2$ & $3$ & $4$ & $5$ & $6$ \\
        $2$ & $2$ & $4$ & $6$ & \cellcolor{seagreen!25}$1$ & $3$ & $5$ \\
        $3$ & $3$ & $6$ & $2$ & $5$ & \cellcolor{seagreen!25}$1$ & $4$ \\
        $4$ & $4$ & \cellcolor{seagreen!25}$1$ & $5$ & $2$ & $6$ & $3$ \\
        $5$ & $5$ & $3$ & \cellcolor{seagreen!25}$1$ & $6$ & $4$ & $2$ \\
        $6$ & $6$ & $5$ & $4$ & $3$ & $2$ & \cellcolor{seagreen!25}$1$ \\
    \end{tabular}
    \end{center}
    In $\mbZ_6$, con $6 = 2 \cdot 3$ non primo, $[2][3] = [0]$: la fattorizzazione di $6$ produce divisori dello zero, e le righe di $2, 3, 4$ hanno
    ripetizioni e non contengono $1$. In $\mbZ_7$ ogni riga è una permutazione di $1, \dots, 6$ e contiene l'$1$: ogni elemento non nullo è
    invertibile ($[3]^{-1} = [5]$, perché $3 \cdot 5 = 15 = 2 \cdot 7 + 1$).
\end{spiegazione}
